\chapter{Módulo \textit{EEG Processor} - Documentación Técnica}
\label{chap:anexo_eeg_processor}

\section{Introducción al Módulo}
\label{sec:intro_eeg_processor}

El módulo \textit{EEG Processor} constituye el núcleo especializado del sistema para captura, procesamiento y análisis de \glspl{biosenial} \acrshort{eeg} en tiempo real. Su desarrollo integra una arquitectura orientada a servicios que soporta múltiples fabricantes de dispositivos comerciales, algoritmos de análisis neurofisiológico avanzados, y capacidades de visualización científica de calidad profesional.

\subsection{Objetivos Técnicos}
\label{subsec:objetivos_eeg_processor}

El desarrollo del módulo se orientó hacia objetivos técnicos específicos para el procesamiento de neuroseñales:

\begin{itemize}
    \item \textbf{Abstracción Multi-Hardware}: Soporte unificado para más de 20 fabricantes mediante \textit{BrainFlow}
    \item \textbf{Procesamiento Tiempo Real}: Algoritmos optimizados para análisis con latencia mínima
    \item \textbf{Análisis Avanzado}: Integración de \acrshort{ica}, transformadas \gls{wavelet} y \textit{MNE-Python}
    \item \textbf{Calidad Científica}: Visualización y exportación orientadas a investigación
    \item \textbf{Arquitectura Extensible}: Servicios modulares para integración con sistemas externos
\end{itemize}

\section{Documentación Técnica Detallada}
\label{sec:doc_tecnica_eeg_processor}

La documentación técnica completa del módulo, incluyendo especificaciones de arquitectura orientada a servicios, algoritmos de procesamiento de señales, integración con dispositivos comerciales, y capacidades de análisis avanzado, se encuentra disponible en el siguiente documento:

\includepdf[pages=-]{imaxes/eeg_processor_doc.pdf}

\section{Visualizadores Especializados}
\label{sec:visualizadores_eeg}

El módulo incluye visualizadores especializados diseñados para diferentes aspectos del análisis de señales \acrshort{eeg}, proporcionando herramientas de visualización científica que facilitan tanto el monitoreo en tiempo real como el análisis detallado de \glspl{biosenial}.

\subsection{Visualizador Multi-Señal}
\label{subsec:visualizador_multi_senal}

El \textit{MultiSignalVisualizer} constituye el componente principal de visualización, proporcionando capacidades de renderizado en tiempo real para múltiples tipos de \glspl{biosenial}. Como se muestra en la Figura~\ref{fig:multi_signal_viz}, el visualizador implementa detección inteligente de tipos de canal, separando automáticamente las señales en pestañas especializadas para \acrshort{eeg}, \acrshort{emg}, \acrshort{ecg}, acelerómetro, giroscopio y otras modalidades de biosensores.

\begin{figure}[hp!]
\centering
\includegraphics[width=\textwidth,height=7cm,keepaspectratio]{imaxes/multi_signal_viz.png}
\caption{Visualizador multi-señal mostrando canales EEG en tiempo real con detección automática de tipos de canal y organización por pestañas especializadas}
\label{fig:multi_signal_viz}
\end{figure}

\subsection{Mapa Topográfico Cerebral}
\label{subsec:mapa_topografico}

El \textit{Brain Map Visualizer} proporciona representación topográfica de la actividad cerebral distribuida sobre el cuero cabelludo, implementando el sistema de posicionamiento estándar 10-20/10-10 para ubicación de electrodos. La Figura~\ref{fig:brain_map_viz} ilustra las capacidades de mapeo en tiempo real con indicadores de calidad de señal codificados por color, permitiendo monitoreo inmediato del estado de conexión de electrodos y distribución espacial de la actividad neuronal.

\begin{figure}[hp!]
\centering
\includegraphics[width=\textwidth,height=7cm,keepaspectratio]{imaxes/brain_map_viz.png}
\caption{Brain Map Visualizer con representación topográfica en tiempo real, mostrando posicionamiento estándar de electrodos y indicadores de calidad de señal por código de colores}
\label{fig:brain_map_viz}
\end{figure}

\subsection{Validación con Hardware Comercial}
\label{subsec:validacion_hardware}

La validación del módulo con dispositivos comerciales demuestra la efectividad de la arquitectura desarrollada. La Figura~\ref{fig:g_tec_suite} presenta la interfaz del \textit{software} propietario \textit{Unicorn Suite} de \textit{g.tec} utilizado con el dispositivo \textit{g.tec Unicorn Hybrid Black}. Esta comparación con las capacidades mostradas en las figuras anteriores evidencia que el módulo desarrollado proporciona funcionalidades de visualización equivalentes a las herramientas comerciales especializadas, manteniendo la flexibilidad de un sistema de código abierto.

\begin{figure}[hp!]
\centering
\includegraphics[width=\textwidth,height=5cm,keepaspectratio]{imaxes/g_tec_suite.png}
\caption{Interfaz de la suite propietaria Unicorn de g.tec utilizada con el dispositivo g.tec Unicorn Hybrid Black, mostrando las capacidades de referencia del software comercial}
\label{fig:g_tec_suite}
\end{figure}

\section{Casos de Uso Principales}
\label{sec:casos_uso_eeg}

\subsection{Investigación Neurocientífica Avanzada}
\label{subsec:investigacion_neurocientifica}

El módulo proporciona herramientas especializadas para investigación con soporte para más de 20 fabricantes de dispositivos comerciales, algoritmos de análisis avanzados incluyendo \acrshort{ica} con 1-16 componentes configurables, integración completa con \textit{MNE-Python} para análisis neurofisiológico, y visualización de calidad científica con exportación de alta resolución para los \gls{ica_analysis}.

\subsection{Desarrollo de Interfaces Cerebro-Computadora}
\label{subsec:desarrollo_bci}

Las capacidades de procesamiento en tiempo real facilitan extracción de características con latencia mínima, cálculo de métricas cognitivas basadas en \gls{machine_learning}, análisis de \gls{band_powers} en tiempo real (Delta, Theta, Alpha, Beta, Gamma), y evaluación continua de calidad de señal.

\subsection{Aplicaciones Clínicas y Diagnósticas}
\label{subsec:aplicaciones_clinicas}

El sistema de monitoreo de calidad de señal en tiempo real, detección automática de \glspl{artefacto} mediante análisis \acrshort{ica}, grabación multi-formato con inserción precisa de marcadores, y compatibilidad con estándares clínicos proporcionan herramientas especializadas para entornos médicos.

\subsection{Análisis de Neuroseñales Multi-Modal}
\label{subsec:analisis_multimodal}

La detección inteligente de tipos de canal (\acrshort{eeg}, \acrshort{emg}, \acrshort{ecg}, EOG, EDA, PPG, acelerómetro, giroscopio), visualización especializada por modalidad de señal, procesamiento coordinado de múltiples fuentes de datos, y análisis correlacional entre diferentes tipos de \glspl{biosenial} amplían las capacidades hacia estudios integrales.